\documentclass[11pt]{article}
\usepackage[margin=1in]{geometry}
\usepackage{amsmath,amssymb,amsthm}
\usepackage{hyperref}
\usepackage{enumitem}

\title{Project Report: Iwasawa Decomposition for \texorpdfstring{$GL(n,\mathbb{R})$}{GL(n,R)} in Lean 4}
\author{LeanCourse25 Project}
\date{\today}

\begin{document}
\maketitle

\begin{abstract}
This project develops a matrix-level Lean 4 framework for the Iwasawa decomposition of
$GL(n,\mathbb{R})$, inspired by Goldfeld's presentation.
The current code formalizes core triangular-matrix infrastructure, a generalized upper-half-space model,
and an API for Iwasawa factorization. The main engineering focus is a clean separation between
foundational lemmas, constructive solver interfaces, and high-level decomposition statements.
\end{abstract}

\section{Goal and Scope}
The target mathematical statement is the decomposition
\[
g = zkd, \qquad g \in GL(n,\mathbb{R}),
\]
with:
\begin{itemize}[leftmargin=1.5em]
\item $z$ in a generalized upper-half-space model $H_n$,
\item $k$ orthogonal,
\item $d$ central scalar (unit in $\mathbb{R}$).
\end{itemize}
The formalization is implemented in
\texttt{LeanCourse25/Projects/Iwasawa.lean}.

\section{Code Organization}
The file is structured into modules:
\begin{itemize}[leftmargin=1.5em]
\item \textbf{Core predicates}: triangular, diagonal, unipotent, diagonal/scalar constructors.
\item \textbf{Triangular lemmas}: closure under addition and multiplication.
\item \textbf{Proof tools}: entrywise extensionality and helper lemmas.
\item \textbf{GL layer}: matrix and invertible-matrix wrappers over $\mathbb{R}$.
\item \textbf{Triangular solver API}: data types and specification for constructive elimination.
\item \textbf{$H_n$ model}: upper-unipotent times normalized positive diagonal.
\item \textbf{Iwasawa API}: factorization record, existence typeclass, quotient target.
\end{itemize}

\section{Architecture Overview (6 Layers)}
The current file follows a strict layered architecture:
\[
\text{Layer 1} \rightarrow \text{Layer 2} \rightarrow \text{Layer 3} \rightarrow
\text{Layer 4} \rightarrow \text{Layer 5} \rightarrow \text{Layer 6}.
\]

\begin{enumerate}[leftmargin=1.5em]
\item \textbf{Layer 1 (L48--262): Matrix predicates and types.}\\
Definitions of \texttt{IsUpperTriangular}, \texttt{IsUpperUnipotent},
\texttt{IsDiagonal}, \texttt{IsOrthogonal}, \texttt{Mat m}, \texttt{GLn m}.
\item \textbf{Layer 2 (L276--714): Triangular solver algorithm.}\\
Algorithmic elimination core:
base case \texttt{solve2\_kill01}, Schur recursion (\texttt{SchurStep}),
and dimension-indexed solver family (\texttt{RecursiveSolveRight}).
\item \textbf{Layer 3 (L748--836): Generalized upper half-plane \texttt{Hn.H m}.}\\
Data model \(z=(x,y)\) with upper-unipotent \(x\), normalized positive diagonal \(y\),
and matrix realization \(\texttt{z.toMat = x * diag(y)}\).
\item \textbf{Layer 4 (L848--901): Factorization data type.}\\
\texttt{IwasawaFactorization} stores \((z,k,d)\) with
\(\texttt{toMat} = z * k * d \cdot I\).
\item \textbf{Layer 5 (L891--1487): Concrete pipeline machinery.}\\
Connects solver outputs to factorization data via
\texttt{ConcretePipeline}, \texttt{FromTriangularSolver}, and \texttt{ConstructiveCore}.
\item \textbf{Layer 6 (L1490--2230): Theorems and \(GL(n,\mathbb{Z})\)-action.}\\
Global existence packaging, uniqueness, quotient equivalence, arithmetic group action,
and Siegel-set statement.
\end{enumerate}

\subsection*{Layer 6 Status Table}
\begin{center}
\begin{tabular}{|l|c|l|}
\hline
\textbf{Theorem} & \textbf{Line} & \textbf{Status} \\
\hline
\texttt{unique\_z\_of\_two\_factorizations} & L1790 & DONE \\
\texttt{Hn\_equiv\_quotient} & L2044 & DONE \\
\texttt{GLZ\_mulAction\_Hn} & L2252 & DONE \\
\texttt{siegel\_fundamental\_domain} & L2279 & \texttt{sorry} \\
\hline
\end{tabular}
\end{center}

\section{Core Formal Definitions}
For $N = \mathrm{Fin}\;m$, the project defines:
\begin{itemize}[leftmargin=1.5em]
\item \texttt{IsUpperTriangular}, \texttt{IsLowerTriangular}, \texttt{IsDiagonal},
\item \texttt{IsUpperUnipotent}, \texttt{IsLowerUnipotent},
\item \texttt{diag : (N -> alpha) -> Matrix N N alpha},
\item \texttt{scalarMat : alpha -> Matrix N N alpha}.
\end{itemize}
These abstractions provide a uniform basis for all later constructions.

\section{Key Proved Lemmas}
The file includes reusable structural lemmas:
\begin{itemize}[leftmargin=1.5em]
\item \texttt{upper\_add}, \texttt{lower\_add},
\item \texttt{upper\_mul}, \texttt{lower\_mul},
\item \texttt{diagonal\_upper}, \texttt{diagonal\_lower},
\item \texttt{upper\_eq\_lower\_implies\_diagonal}.
\end{itemize}
These are the main proof-theoretic building blocks for uniqueness arguments and elimination workflows.

\section{Constructive Solver Layer}
In namespace \texttt{TriangularSolver}, the code introduces:
\begin{itemize}[leftmargin=1.5em]
\item \texttt{UpperUnipotent} as a structure carrying both data and proofs,
\item \texttt{SolveRightSpec} as a solver contract for equations of the form $U A = B$,
\item completeness formalization \texttt{SolveRightSpec.Complete} and impossibility result \texttt{not\_complete},
\item concrete $2\times 2$ lemmas \texttt{solve2\_kill01}, \texttt{solve2\_upperUnipotent}, \texttt{solve2\_lowering\_data},
\item a field-level solver \texttt{solveRight2Field} and its real specialization \texttt{solveRight2Real},
\item recursive infrastructure \texttt{SchurStep}, \texttt{stepOfSchur}, \texttt{recursiveSolveRightOfSchur}, and \texttt{RecursiveSolveRight}.
\end{itemize}
This design isolates computational elimination from downstream decomposition logic.

\section{Generalized Upper Half-Space \texorpdfstring{$H_n$}{Hn}}
The file models $H_n$ as structured data:
\[
z = x y,
\]
where:
\begin{itemize}[leftmargin=1.5em]
\item $x$ is upper-unipotent (\texttt{Hn.UpperUnipotent}),
\item $y$ is normalized positive diagonal (\texttt{Hn.PosDiagNorm}).
\end{itemize}
The map \texttt{Hn.H.toMat} gives the underlying matrix representation.

\section{Iwasawa API}
The final API layer defines:
\begin{itemize}[leftmargin=1.5em]
\item \texttt{IwasawaFactorization} with fields $(z,k,d)$,
\item \texttt{ConcretePipeline} (algorithmic data: \texttt{xOf}, \texttt{yOf}, \texttt{qOf}, \texttt{dOf}, \texttt{correct}),
\item bridge layers \texttt{FromTriangularSolver} and \texttt{ConstructiveCore},
\item \texttt{HasIwasawa} (typeclass) and immediate theorem \texttt{exists\_factorization},
\item uniqueness theorem \texttt{unique\_z\_of\_two\_factorizations},
\item quotient objects \texttt{OZ\_subgroup}, \texttt{Hquot}, \texttt{Hn\_equiv\_quotient}.
\end{itemize}
This makes downstream usage independent of the eventual constructive proof of the instance.

\subsection*{Core Contract: \texttt{HasIwasawa}}
The central abstraction is:
\begin{itemize}[leftmargin=1.5em]
\item a constructor
\[
\texttt{iwasawa : GLn m -> IwasawaFactorization m},
\]
which assigns factors $(z,k,d)$ to each invertible matrix $g$;
\item a correctness theorem
\[
\texttt{correct : }\forall g,\; (\texttt{iwasawa}\;g).\texttt{toMat} = g.
\]
\end{itemize}

So \texttt{HasIwasawa} is exactly ``existence + reconstruction correctness''.
All downstream theorems can depend on this interface without committing to a specific algorithm.

\subsection*{Why \texttt{exists\_factorization} is immediate}
The proof
\[
\texttt{theorem exists\_factorization [HasIwasawa m] (g : GLn m) :}
\]
\[
\texttt{  exists fac, fac.toMat = g}
\]
is one-step logical unpacking of the class fields:
\begin{enumerate}[leftmargin=1.5em]
\item choose the witness
\[
\texttt{fac := HasIwasawa.iwasawa g};
\]
\item remaining goal is exactly
\[
\texttt{fac.toMat = g},
\]
which is \texttt{HasIwasawa.correct g};
\item \texttt{simpa} finishes after unfolding the chosen witness.
\end{enumerate}
In short, the theorem is a canonical projection from the typeclass data.
\subsection*{Full Construction Pipeline (Exact Code Path)}
Below is the concrete path in your code from input \(g\) to the data used by
\texttt{HasIwasawa}.
\begin{enumerate}[leftmargin=1.5em]
\item \textbf{Input.}
Start from
\[
g : \texttt{GLn m}.
\]
Target:
\[
\texttt{fac : IwasawaFactorization m}
\]
\item \textbf{Base elimination at dimension 2.}
Use \texttt{solve2\_kill01} to kill entry \((0,1)\), then package as
\texttt{UpperUnipotent} via \texttt{solve2\_upperUnipotent}, and package
\(B=U*A\) via \texttt{solve2\_lowering\_data}. This yields the concrete base
\texttt{solveRight2Field} and \texttt{solveRight2Real}.

\item \textbf{Recursive solver family.}
Provide Schur-step data \texttt{SchurStep} and build higher-dimensional solvers
with \texttt{stepOfSchur} and \texttt{recursiveSolveRightOfSchur}, giving
\texttt{RecursiveSolveRight}.

\item \textbf{Algorithmic factor data at fixed dimension.}
Use \texttt{ConcretePipeline m}, whose fields are exactly the data production
functions:
\texttt{xOf}, \texttt{yOf}, \texttt{qOf}, \texttt{qOrth}, \texttt{dOf},
plus a reconstruction proof \texttt{correct}.

\item \textbf{Package factors.}
From any pipeline \texttt{P : ConcretePipeline m} and \(g\):
\begin{itemize}[leftmargin=1.5em]
\item \texttt{zOf P g : Hn.H m},
\item \texttt{facOf P g : IwasawaFactorization m},
\item \texttt{facOf\_correct P g : (facOf P g).toMat = g}.
\end{itemize}
This already gives \texttt{exists\_factorization\_of\_pipeline}.

\item \textbf{Dimension recursion for pipelines.}
Use \texttt{RecursiveMkPipeline} and \texttt{RecursiveConcretePipeline} to get
pipelines at arbitrary dimensions \(m \ge 2\), with theorem
\texttt{RecursiveConcretePipeline.exists\_factorization\_ge2}.

\item \textbf{Bridge from solver recursion to pipeline recursion.}
Use \texttt{FromTriangularSolver} (fields: \texttt{solver}, \texttt{mkPipeline}),
or equivalently raw witness packaging via \texttt{RawMkPipeline}.
This yields theorems
\texttt{exists\_raw\_ge2}, \texttt{exists\_factorization\_ge2},
and \texttt{Fact}-versions.

\item \textbf{Single constructive package.}
Bundle everything into
\[
\texttt{ConstructiveCore := \{ solverBase, solverSchur, pipelineBuilder \}}.
\]
Then derive:
\texttt{solverRec}, \texttt{bridge},
\texttt{exists\_factorization\_fact}, and \texttt{toHasIwasawaFact}.

\item \textbf{Typeclass extraction.}
At this point, obtain
\[
\texttt{HasIwasawa m}
\]
via \texttt{ConcretePipeline.toHasIwasawa},
\texttt{FromTriangularSolver.toHasIwasawaFact}, or
\texttt{ConstructiveCore.toHasIwasawaFact}.
With a registered global core (\texttt{HasConstructiveCore}), the instance
\texttt{hasIwasawaFact} is automatic.

\item \textbf{Final theorem as projection.}
Once \([\texttt{HasIwasawa m}]\) is available, the theorem
\texttt{exists\_factorization} is immediate by taking witness
\texttt{HasIwasawa.iwasawa g} and applying \texttt{HasIwasawa.correct g}.
\end{enumerate}

\noindent
Notation map used in this report:
\[
(U,M,D)\ \longrightarrow\ (x,y,d)\ \longrightarrow\ (z,k,d),
\]
where \(z=(x,y)\) is exactly the \texttt{Hn.H}-object stored in
\texttt{IwasawaFactorization}.

\section{Current Status and Remaining Tasks}
The infrastructure and API are in place and compile.
Major completed milestones now include:
\begin{itemize}[leftmargin=1.5em]
\item constructive base solver and recursive solver framework (\texttt{TriangularSolver}),
\item concrete and recursive pipeline packaging (\texttt{ConcretePipeline}, \texttt{ConstructiveCore}),
\item uniqueness of the \(z\)-factor (\texttt{unique\_z\_of\_two\_factorizations}),
\item quotient construction and equivalence map (\texttt{OZ\_subgroup}, \texttt{Hquot}, \texttt{Hn\_equiv\_quotient}).
\end{itemize}
At the moment, the remaining unfinished theorem in this file is:
\begin{itemize}[leftmargin=1.5em]
\item \texttt{siegel\_fundamental\_domain} (currently marked with \texttt{sorry}).
\end{itemize}

\section{Conclusion}
This project establishes a strong formal skeleton for Iwasawa decomposition in Lean 4:
core matrix algebra, triangular proof infrastructure, a solver interface, and a scalable API.
The architecture is suitable for extending from specification-level statements to full constructive proofs.

\section*{Reference}
D. Goldfeld, \emph{Automorphic Forms and L-Functions for the Group GL(n, R)},
Cambridge Studies in Advanced Mathematics, vol.~99, 2006.

\end{document}
